%% paper81_hadwiger_vermutung_de.tex
%% Hadwigers Vermutung für n >= 7: Graphenfärbung und Minortheorie (Deutsch)
%% Autor: Michael Fuhrmann
%% Specialist Mathematics Project, Build 166
%% Datum: 2026-03-12

\documentclass[12pt,a4paper]{amsart}
\usepackage{amsmath,amssymb,amsthm,mathtools,enumitem,booktabs,array,hyperref}
\usepackage[utf8]{inputenc}
\usepackage[T1]{fontenc}
\usepackage[ngerman]{babel}
\usepackage{geometry}
\geometry{margin=2.5cm}

%% Theorem-Umgebungen (Deutsch)
\newtheorem{theorem}{Satz}[section]
\newtheorem{satz}[theorem]{Satz}
\newtheorem{lemma}[theorem]{Lemma}
\newtheorem{korollar}[theorem]{Korollar}
\newtheorem{corollary}[theorem]{Korollar}
\newtheorem{proposition}[theorem]{Proposition}
\newtheorem{vermutung}[theorem]{Vermutung}
\newtheorem{conjecture}[theorem]{Vermutung}
\theoremstyle{definition}
\newtheorem{definition}[theorem]{Definition}
\newtheorem{definition_de}[theorem]{Definition}
\newtheorem{beispiel}[theorem]{Beispiel}
\theoremstyle{remark}
\newtheorem{bemerkung}[theorem]{Bemerkung}
\newtheorem{remark}[theorem]{Bemerkung}

\hypersetup{
  colorlinks=true,
  linkcolor=blue,
  citecolor=blue,
  urlcolor=blue
}

\begin{document}

\title{Hadwigers Vermutung für $k \geq 7$:\\
Graphenfärbung, vollständige Minoren\\
und die Verbindung zum Vier-Farben-Satz}

\author{Michael Fuhrmann}
\address{Specialist Mathematics Project, Build 166}
\email{specialist-maths@project.local}
\date{2026-03-12}

\begin{abstract}
Hadwigers Vermutung (1943) behauptet, dass jeder Graph $G$ mit chromatischer
Zahl $\chi(G) \geq k$ den vollständigen Graphen $K_k$ als Minor enthält. Dies
ist eines der tiefsten offenen Probleme der Graphentheorie, das die Theorie der
Graphenfärbungen mit der Strukturtheorie der Graphenminoren nach Robertson und
Seymour verbindet.

Die Vermutung ist für $k \leq 6$ verifiziert: Die Fälle $k \leq 4$ folgen aus
klassischen Resultaten, $k = 5$ ist äquivalent zum Vier-Farben-Satz (Wagner 1937),
und $k = 6$ wurde von Robertson, Seymour und Thomas (1993) ebenfalls unter Nutzung
des Vier-Farben-Satzes bewiesen. Für $k \geq 7$ ist die Vermutung vollständig offen.

Wir präsentieren vollständige Beweise für $k \leq 4$, eine detaillierte Analyse
der $k = 5$-Äquivalenz mit dem Vier-Farben-Satz, den Robertson--Seymour--Thomas-Beweis
für $k = 6$ sowie Teilresultate für große $k$, einschließlich der Reed--Seymour-Schranken
und der Ergebnisse von Kühn--Osthus.
\end{abstract}

\maketitle

\tableofcontents

%% =============================================================================
\section{Einleitung}
%% =============================================================================

Im Jahr 1943 stellte Hugo Hadwiger~\cite{Hadwiger1943} folgende vereinigende Vermutung vor,
die die chromatische Zahl eines Graphen mit seiner Minorstruktur verbindet:

\begin{vermutung}[Hadwiger, 1943]\label{verm:hadwiger}
Für jede ganze Zahl $k \geq 1$ gilt: Falls $\chi(G) \geq k$, dann enthält $G$
den vollständigen Graphen $K_k$ als Minor ($G \succeq K_k$).
\end{vermutung}

Äquivalent: Jeder Graph ohne $K_k$-Minor ist $(k-1)$-färbbar.

\begin{definition_de}
Sei $G = (V,E)$ ein Graph.
\begin{enumerate}[label=(\alph*)]
  \item Die \emph{chromatische Zahl} $\chi(G)$ ist die kleinste Anzahl an Farben,
        die benötigt wird, um die Knoten von $G$ so zu färben, dass benachbarte
        Knoten verschiedene Farben erhalten.
  \item Ein Graph $H$ ist ein \emph{Minor} von $G$ (geschrieben $H \preceq G$),
        falls $H$ aus $G$ durch eine Folge von Kantenkon\-trak\-tionen, Kantenlöschungen
        und Knotenlöschungen gewonnen werden kann.
  \item Die \emph{Hadwiger-Zahl} $\eta(G)$ (auch $h(G)$) ist das größte $k$
        mit $K_k \preceq G$.
\end{enumerate}
\end{definition_de}

Hadwigers Vermutung lautet umformuliert: $\chi(G) \leq \eta(G)$ für alle Graphen $G$.

%% =============================================================================
\section{Graphenminoren: Grundlegende Theorie}
%% =============================================================================

\begin{definition_de}
Eine \emph{Kantenkontraktion} kontrahiert eine Kante $uv$, indem $u$ und $v$ zu
einem einzelnen Knoten $w$ mit Nachbarschaft $N(u) \cup N(v) \setminus \{u,v\}$
zusammengeführt werden.
\end{definition_de}

\begin{proposition}\label{prop:minoreig}
Die Minorrelation $\preceq$ ist eine Halbordnung auf Isomorphieklassen von Graphen.
Falls $H \preceq G$, so gilt $\chi(H) \leq \chi(G)$.
\end{proposition}

\begin{satz}[Robertson--Seymour, Graphenminor-Satz 2004]\label{satz:GMT}
In jeder unendlichen Folge von Graphen $G_1, G_2, \ldots$ gibt es zwei Graphen
$G_i \preceq G_j$ mit $i < j$. Insbesondere hat jede minorabgeschlossene Graphenfamilie
eine endliche Menge verbotener Minoren.
\end{satz}

%% =============================================================================
\section{Bewiesene Fälle: \texorpdfstring{$k \leq 4$}{k <= 4}}
%% =============================================================================

\begin{satz}[Hadwigers Vermutung für $k \leq 3$]\label{satz:k3}
\begin{enumerate}[label=(\roman*)]
  \item $k = 1$: Jeder nichtleere Graph hat $\chi(G) \geq 1$ und enthält $K_1$.
  \item $k = 2$: $\chi(G) \geq 2$ impliziert, $G$ hat eine Kante, also $K_2 \preceq G$.
  \item $k = 3$: $\chi(G) \geq 3$ impliziert, $G$ enthält einen ungeraden Kreis, also $K_3 \preceq G$.
\end{enumerate}
\end{satz}

\begin{proof}
Für $k = 3$: Da $\chi(G) \geq 3$, ist $G$ nicht bipartit, enthält also einen ungeraden
Kreis $C_{2m+1}$, $m \geq 1$. Jeder ungerade Kreis enthält $K_3$ als Minor
(durch Kontrahieren aller bis auf drei Knoten), also $K_3 \preceq G$.
\end{proof}

\begin{satz}[Hadwiger für $k = 4$, Dirac 1952]\label{satz:k4}
Jeder $4$-chromatische Graph enthält $K_4$ als Minor.
\end{satz}

\begin{proof}[Beweisidee]
Dirac~\cite{Dirac52} bewies, dass jeder $4$-chromatische Graph $K_4$ als topologischen
Minor enthält (d.h.\ eine Unterteilung von $K_4$). Da topologischer Minor ein Spezialfall
von Minor ist, folgt das Resultat. Der Beweis nutzt Kempe-Ketten und die Nichtplanarität
$4$-chromatischer Graphen.
\end{proof}

%% =============================================================================
\section{Wagners Satz und der Fall \texorpdfstring{$k = 5$}{k=5}}
%% =============================================================================

\begin{satz}[Wagner, 1937]\label{satz:wagner}
Ein Graph $G$ ist genau dann planar, wenn er weder $K_5$ noch $K_{3,3}$ als Minor enthält.
\end{satz}

\begin{satz}[Vier-Farben-Satz, Appel--Haken 1977; Robertson et al.\ 1997]\label{satz:VFS}
Jeder planare Graph ist $4$-färbbar.
\end{satz}

\begin{satz}[Hadwiger für $k = 5$ $\Leftrightarrow$ VFS, Wagner 1937]\label{satz:k5}
Hadwigers Vermutung für $k = 5$ ist äquivalent zum Vier-Farben-Satz.
\end{satz}

\begin{proof}
($\Leftarrow$) Angenommen, der Vier-Farben-Satz gelte. Sei $G$ ein Graph mit
$\chi(G) \geq 5$. Dann ist $G$ nach dem VFS nicht planar. Nach Wagners Satz enthält
$G$ dann $K_5$ oder $K_{3,3}$ als Minor. Mithilfe von Wagners Struktursatz
(Graphen ohne $K_5$-Minor sind aus planaren Graphen und dem Wagner-Graphen $V_8$
durch $3$-Summen aufgebaut) folgt, dass $K_{3,3}$-Minor-freie $5$-chromatische Graphen
stets $K_5$ als Minor enthalten. Also $K_5 \preceq G$.

($\Rightarrow$) Gilt Hadwigers Vermutung für $k = 5$, so hat jeder planare Graph
(der keinen $K_5$-Minor enthält) chromatische Zahl $\leq 4$.
\end{proof}

%% =============================================================================
\section{Der Fall \texorpdfstring{$k = 6$}{k=6}: Robertson--Seymour--Thomas}
%% =============================================================================

\begin{satz}[Robertson--Seymour--Thomas, 1993]\label{satz:k6}
Jeder $6$-chromatische Graph enthält $K_6$ als Minor. Äquivalent: Jeder Graph
ohne $K_6$-Minor ist $5$-färbbar.
\end{satz}

\begin{proof}[Beweisübersicht]
Robertson, Seymour und Thomas~\cite{RST93} zeigen, dass jeder minimale Gegenbeispiel-Kandidat
zu Hadwigers Vermutung für $k = 6$ einen \emph{Apex-Knoten} $v$ besitzen muss, sodass
$G - v$ planar ist. Da $G - v$ planar ist, gilt nach dem VFS $\chi(G - v) \leq 4$.
Durch Zuweisen einer fünften Farbe an $v$ ergibt sich $\chi(G) \leq 5$, ein Widerspruch
zur Annahme $\chi(G) = 6$. Die Strukturzerlegung von $K_6$-minor-freien Graphen
ist das Herzstück des Beweises.
\end{proof}

\begin{korollar}
Die Fälle $k \leq 6$ von Hadwigers Vermutung sind bewiesene Sätze. Für $k = 5, 6$
benötigen die Beweise den Vier-Farben-Satz.
\end{korollar}

%% =============================================================================
\section{Minorabgeschlossene Graphenfamilien und Wagners Formel}
%% =============================================================================

\begin{satz}[Kostochka 1984; Thomason 1984]\label{satz:kostochka}
Es gibt eine Konstante $c > 0$, sodass jeder Graph $G$ mit mittlerem Grad
$\geq c \cdot k \sqrt{\log k}$ den Graphen $K_k$ als Minor enthält. Die optimale
Konstante ist $c = (\alpha + o(1))/\sqrt{\ln k}$ mit $\alpha = 0{,}319\ldots$\,.
\end{satz}

\begin{satz}[Mader, 1967]\label{satz:mader}
Jeder Graph ohne $K_k$-Minor hat mittleren Grad $< 2^{k-2}$. Insbesondere ist jeder
$K_k$-minorfreie Graph $(2^{k-2})$-degeneriert, also $(2^{k-2}+1)$-färbbar.
\end{satz}

\begin{bemerkung}
Die exponentielle Lücke zwischen Maders Schranke $2^{k-2}$ und Hadwigers Vorhersage
$k-1$ ist für große $k$ enorm. Für $k = 5$: Mader liefert $9$-Färbbarkeit (da $2^{5-2}+1 = 9$) ohne VFS,
für $k = 6$: nur $17$-Färbbarkeit (da $2^{6-2}+1 = 17$), aber RST verbessert auf $5$.
\end{bemerkung}

%% =============================================================================
\section{Kühn--Osthus und große chromatische Zahlen}
%% =============================================================================

\begin{satz}[Kühn--Osthus, 2005]\label{satz:KO}
Für jedes $\varepsilon > 0$ gibt es $k_0$, sodass für alle $k \geq k_0$ gilt:
Falls $\chi(G) \geq k$, dann $\eta(G) \geq k/(1+\varepsilon)$.
\end{satz}

\begin{proof}[Beweisidee]
Jeder $k$-chromatische Graph enthält einen $k$-kritischen Teilgraphen mit
Minimalgrad $\geq k-1$. Der Thomason-Satz verbindet hohen Minimalgrad mit
großer Hadwiger-Zahl; der Faktor $(1+\varepsilon)$ kompensiert die Lücke
zwischen Minimal- und Durchschnittsgrad.
\end{proof}

\begin{satz}[Norin--Song, 2023]\label{satz:NS23}
Für alle hinreichend großen $k$ gilt: Falls $\chi(G) \geq k$, dann
$\eta(G) \geq k/(\log k)^{0{,}999}$.
\end{satz}

Dies ist das derzeit beste asymptotische Resultat, liegt aber noch weit von der
linearen Schranke $\eta(G) \geq k$ entfernt, die Hadwigers Vermutung fordert.

%% =============================================================================
\section{Gebrochene und lineare Hadwiger-Vermutung}
%% =============================================================================

\begin{definition_de}
Die \emph{gebrochene chromatische Zahl} $\chi_f(G)$ ist das Infimum über $p/q$,
für das eine zulässige $p$-Färbung von $G^q$ (dem $q$-fachen lexikographischen
Potenz von $G$) existiert. Es gilt $\chi_f(G) \leq \chi(G)$.
\end{definition_de}

\begin{satz}[Gebrochener Hadwiger, Reed--Seymour 1998]\label{satz:fracHad}
$\chi_f(G) \leq 2\eta(G)$ für alle Graphen $G$.
\end{satz}

\begin{vermutung}[Linearer Hadwiger]\label{verm:linHad}
Es gibt eine absolute Konstante $c$, sodass $\chi(G) \leq c \cdot \eta(G)$
für alle Graphen $G$.
\end{vermutung}

\begin{bemerkung}
Hadwigers Vermutung ist der Spezialfall $c = 1$. Die lineare Hadwiger-Vermutung
fordert irgendein festes $c$ und wäre bereits ein bedeutender Fortschritt.
Das beste aktuelle Ergebnis (Norin--Song 2023) liefert $c = O((\log k)^{0{,}999})$ –
noch logarithmisch.
\end{bemerkung}

%% =============================================================================
\section{Die Ungerade Hadwiger-Vermutung (Gerards--Seymour)}
%% =============================================================================

\begin{vermutung}[Ungerade Hadwiger, Gerards--Seymour]\label{verm:odd}
Falls $G$ keinen $K_k$ als \emph{ungeraden Minor} enthält (d.\,h.\ einen Minor,
bei dem alle Verzweigungsmengen ungerade Größe haben oder eine bestimmte
Paritätsbedingung erfüllen), dann gilt $\chi(G) \leq 2(k-1)$.
\end{vermutung}

\begin{satz}[Geelen, Gerards, Reed, Seymour, Vetta, 2009]\label{satz:oddHad}
Die ungerade Hadwiger-Vermutung gilt für $k \leq 4$.
\end{satz}

%% =============================================================================
\section{Aktueller Stand und offene Probleme}
%% =============================================================================

\begin{vermutung}[Hadwiger für $k \geq 7$ (Offen)]\label{verm:offen7}
Für $k \geq 7$: Jeder Graph mit $\chi(G) \geq k$ enthält $K_k$ als Minor.
\end{vermutung}

\begin{center}
\begin{tabular}{cll}
\toprule
$k$ & Status & Referenz \\
\midrule
$1, 2, 3$ & Trivial & Klassisch \\
$4$ & Bewiesen & Dirac (1952) \\
$5$ & Bewiesen $\Leftrightarrow$ VFS & Wagner (1937) + Appel--Haken (1977) \\
$6$ & Bewiesen (mit VFS) & Robertson--Seymour--Thomas (1993) \\
$7$ & \textbf{Offen} & -- \\
$\geq 7$ & \textbf{Offen} & -- \\
$\geq k_0$ (groß) & $\eta(G) \geq k/(\log k)^{0{,}999}$ & Norin--Song (2023) \\
\bottomrule
\end{tabular}
\end{center}

\medskip

Hauptoffene Fragen:
\begin{enumerate}[label=(\arabic*)]
  \item Beweise oder widerlege Hadwigers Vermutung für $k = 7$.
  \item Gibt es einen Beweis für alle $k$, der den VFS nicht verwendet?
  \item Konstruiere einen Graphen mit $\chi(G) = k$ aber $\eta(G) < k$ (Gegenbeispiel)
        oder beweise, dass keiner existiert.
  \item Verbessere die asymptotische Schranke: zeige $\eta(G) \geq c \cdot \chi(G)$
        für eine Konstante $c > 0$ unabhängig von $\chi(G)$.
  \item Beweise die lineare Hadwiger-Vermutung.
\end{enumerate}

%% =============================================================================
\section{Zusammenfassung}
%% =============================================================================

Hadwigers Vermutung vereint zwei zentrale Themen der Graphentheorie: chromatische
Zahlen und Graphenminoren. Ihr Beweis für $k \leq 6$ erforderte bereits einen der
tiefsten Sätze der Mathematik (den Vier-Farben-Satz), und die Fälle $k \geq 7$
werden als dramatisch schwieriger erwartet.

Die Robertson--Seymour-Theorie der Graphenminoren liefert mächtige Strukturwerkzeuge,
aber selbst diese reichen für $k \geq 7$ nicht aus. Das Problem liegt im Schnittpunkt
extremaler Graphentheorie, topologischer Graphentheorie und Berechnungskomplexität;
seine Lösung würde fundamental neue Ideen erfordern.

\begin{thebibliography}{99}

\bibitem{Hadwiger1943}
H.~Hadwiger,
\textit{Über eine Klassifikation der Streckenkomplexe},
Vierteljschr.\ Naturforsch.\ Ges.\ Zürich \textbf{88} (1943), 133--142.

\bibitem{Wagner37}
K.~Wagner,
\textit{Über eine Eigenschaft der ebenen Komplexe},
Math.\ Ann.\ \textbf{114} (1937), 570--590.

\bibitem{Dirac52}
G.~A.\ Dirac,
\textit{A property of $4$-chromatic graphs and some remarks on critical graphs},
J.\ London Math.\ Soc.\ \textbf{27} (1952), 85--92.

\bibitem{RST93}
N.~Robertson, P.~D.\ Seymour, R.~Thomas,
\textit{Hadwiger's conjecture for $K_6$-free graphs},
Combinatorica \textbf{13} (1993), Nr.\ 3, 279--361.

\bibitem{AppHak77}
K.~Appel, W.~Haken,
\textit{Every planar map is four colorable. Part I: Discharging},
Illinois J.\ Math.\ \textbf{21} (1977), Nr.\ 3, 429--490.
% Ankündigung: Bull.\ Amer.\ Math.\ Soc.\ \textbf{82} (1976), 711--712.

\bibitem{RSST97}
N.~Robertson, D.~Sanders, P.~D.\ Seymour, R.~Thomas,
\textit{The four-colour theorem},
J.\ Combin.\ Theory Ser.\ B \textbf{70} (1997), Nr.\ 1, 2--44.

\bibitem{Kostochka84}
A.~W.\ Kostochka,
\textit{Lower bound of the Hadwiger number of graphs by their average degree},
Combinatorica \textbf{4} (1984), Nr.\ 4, 307--316.

\bibitem{Thomason84}
A.~Thomason,
\textit{An extremal function for contractions of graphs},
Math.\ Proc.\ Cambridge Philos.\ Soc.\ \textbf{95} (1984), 261--265.

\bibitem{KO05}
D.~Kühn, D.~Osthus,
\textit{Forcing unbalanced complete bipartite minors},
European J.\ Combin.\ \textbf{26} (2005), 75--81.

\bibitem{ReedSey98}
B.~Reed, P.~D.\ Seymour,
\textit{Fractional colouring and Hadwiger's conjecture},
J.\ Combin.\ Theory Ser.\ B \textbf{74} (1998), 147--152.

\bibitem{NorinSong23}
S.~Norin, L.~Song,
\textit{Breaking the degeneracy barrier for coloring graphs with no $K_t$ minor},
Adv.\ Math.\ \textbf{422} (2023), Paper Nr.\ 109020.

\bibitem{Mader67}
W.~Mader,
\textit{Homomorphieeigenschaften und mittlere Kantendichte von Graphen},
Math.\ Ann.\ \textbf{174} (1967), 265--268.

\end{thebibliography}

\end{document}
